%% ============================================================
%%  Chapitre 4 — AutoML pour le scoring scientifique
%% ============================================================
\chapter{AutoML pour le scoring scientifique}
\label{ch:automl}

\section{Motivation : les limites des pondérations expertes}

La combinaison linéaire des six dimensions de qualité définie à
l'équation~\ref{eq:weighted_score} soulève une question fondamentale :
comment choisir les pondérations $w_f$ ? Une approche naïve consiste
à fixer des pondérations expertes uniformes ($w_f = 1/6 \approx 0.167$)
ou à les calibrer manuellement en se fondant sur l'intuition du domaine.

Cette approche présente deux limitations majeures. D'une part, elle suppose
que les six dimensions contribuent de manière égale à la qualité perçue,
ce qui est inexact : l'importance relative de la vélocité (pour les domaines
d'actualité) versus la récence (pour les domaines plus matures) varie
significativement selon le contexte. D'autre part, une calibration manuelle
est subjective et difficilement reproductible d'un domaine à l'autre.

L'optimisation automatique des pondérations par un algorithme AutoML
constitue une alternative rigoureuse, reproductible et potentiellement
plus performante. C'est l'objet de ce chapitre.

\section{Formalisation du problème d'optimisation}

\subsection{Espace de recherche}

L'espace de recherche est le simplexe unitaire en dimension 6 :

\begin{equation}
  \mathcal{W} = \left\{ \mathbf{w} \in \mathbb{R}^6 \mid
    w_f \geq 0 \;\forall f,\; \sum_{f} w_f = 1 \right\}
  \label{eq:simplex}
\end{equation}

En pratique, Optuna optimise les six pondérations dans $[0, 1]$
indépendamment, puis normalise le vecteur résultant pour satisfaire
la contrainte de somme unitaire. Ce reparamétrisation préserve
la topologie de l'espace de recherche tout en simplifiant l'exploration.

\subsection{Fonction objectif}

La fonction objectif est le NDCG@15 calculé sur l'oracle du domaine :

\begin{equation}
  \mathbf{w}^* = \arg\max_{\mathbf{w} \in \mathcal{W}}
    \text{NDCG@15}\!\left(\text{rank}_{\mathbf{w}}(C_d), G_d\right)
  \label{eq:optimization_objective}
\end{equation}

où $C_d$ est le corpus du domaine $d$, $G_d$ est l'oracle de pertinence
(grades 0, 1, 2 pour chaque article), et $\text{rank}_{\mathbf{w}}(C_d)$
est le classement des articles du corpus selon le score
$\text{score}_{\mathbf{w}}(a) = \sum_f w_f \cdot s_f(a)$.

\subsection{Algorithme : TPE d'Optuna}

Nous utilisons l'échantillonneur TPE (\textit{Tree-structured Parzen
Estimator}) d'Optuna \cite{akiba2019optuna}, configuré avec 150 essais
(\textit{trials}) par domaine. TPE maintient deux distributions
de probabilité : $l(\mathbf{w})$ modélisant les pondérations associées
aux meilleurs essais, et $g(\mathbf{w})$ celles associées aux essais
moins performants. À chaque itération, il sélectionne les pondérations
qui maximisent le ratio $\text{EI}(\mathbf{w}) \propto l(\mathbf{w})/g(\mathbf{w})$,
focalisant progressivement l'exploration sur les régions prometteuses
de $\mathcal{W}$.

\section{Résultats : validation de H1}

\subsection{Amélioration par domaine}

Le tableau~\ref{tab:automl_results} présente les résultats de l'optimisation
AutoML pour les 19 domaines d'évaluation. Pour chaque domaine, nous
reportons le NDCG@15 obtenu avec les pondérations par défaut (uniformes),
les pondérations apprises par Optuna (V2, avec pertinence sémantique),
et l'amélioration relative.

\begin{table}[htbp]
  \centering
  \caption{Résultats AutoML par domaine : NDCG@15 et amélioration relative.}
  \label{tab:automl_results}
  \rowcolors{2}{lightgray}{white}
  \begin{tabular}{lccr}
    \toprule
    \textbf{Domaine} & \textbf{Default} & \textbf{Optuna V2} & \textbf{Amélioration} \\
    \midrule
    fake\_news\_detection  & 0.093 & 0.470 & $+406\,\%$ \\
    neural\_architecture\_search & 0.359 & 0.590 & $+64\,\%$ \\
    knowledge\_graph\_embedding & 0.332 & 0.602 & $+81\,\%$ \\
    graph\_attention\_networks & 0.396 & 0.645 & $+63\,\%$ \\
    explainable\_ai & 0.067 & 0.206 & $+208\,\%$ \\
    text\_summarization & 0.391 & 0.681 & $+74\,\%$ \\
    computational\_biology & 0.470 & 0.566 & $+20\,\%$ \\
    sentiment\_analysis & 0.151 & 0.351 & $+133\,\%$ \\
    deep\_reinforcement\_learning & 0.124 & 0.251 & $+102\,\%$ \\
    \textit{Moyenne (19 domaines)} & \textit{0.229} & \textit{0.387} & \textit{$+68.7\,\%$} \\
    \bottomrule
  \end{tabular}
  \small{\textit{Note : seuls 9 domaines représentatifs sont listés par souci de lisibilité.
  Les 19 résultats complets sont disponibles dans \texttt{data/metalearning/phase8c\_results.json}.}}
\end{table}

\begin{resultbox}
  \textbf{H1 confirmée.} Sur les 19 domaines évalués, l'optimisation Optuna
  améliore systématiquement le NDCG@15 par rapport aux pondérations par défaut,
  avec une amélioration moyenne de $+68.7\,\%$ ($\text{min} = +20\,\%$,
  $\text{max} = +406\,\%$). L'amélioration minimale dépasse largement le
  seuil de $5\,\%$ fixé par H1, confirmant l'hypothèse.
\end{resultbox}

\subsection{Analyse des profils de pondération}

Au-delà des performances agrégées, l'analyse des vecteurs de pondérations
appris révèle des patterns intéressants qui constituent le signal brut
pour l'apprentissage méta (Chapitre~\ref{ch:metalearning}).

Le tableau~\ref{tab:weight_profiles} illustre la diversité des profils
appris sur quelques domaines représentatifs.

\begin{table}[htbp]
  \centering
  \caption{Profils de pondération appris par Optuna (pondérations normalisées).}
  \label{tab:weight_profiles}
  \rowcolors{2}{lightgray}{white}
  \begin{tabular}{lcccccc}
    \toprule
    \textbf{Domaine} & \textbf{venue} & \textbf{authors} & \textbf{impact}
      & \textbf{velocity} & \textbf{recency} & \textbf{relevance} \\
    \midrule
    fake\_news & 0.027 & 0.061 & 0.037 & \textbf{0.698} & 0.104 & 0.024 \\
    federated\_learning & 0.201 & \textbf{0.295} & 0.175 & 0.138 & 0.153 & 0.038 \\
    graph\_NN & 0.204 & 0.128 & 0.084 & 0.221 & \textbf{0.329} & 0.034 \\
    explainable\_ai & 0.089 & 0.074 & 0.214 & \textbf{0.296} & 0.029 & \textbf{0.298} \\
    comp. biology & 0.112 & \textbf{0.281} & 0.168 & 0.165 & 0.210 & 0.063 \\
    quant. finance & \textbf{0.283} & 0.150 & 0.069 & 0.114 & 0.032 & \textbf{0.352} \\
    \bottomrule
  \end{tabular}
\end{table}

Plusieurs observations se dégagent. La pondération de la \textbf{vélocité}
est prépondérante pour la détection de fausses nouvelles (0.698), un domaine
d'actualité où les articles récemment cités sont les plus pertinents.
À l'inverse, pour les \textbf{réseaux de neurones sur graphes} (domaine
plus mature), la récence prend le dessus (0.329). La \textbf{pertinence
sémantique} est fortement valorisée pour l'IA explicable et la finance
quantitative, deux domaines au vocabulaire technique très spécifique.

Ces divergences de profils constituent la manifestation concrète
du phénomène que le méta-learning cherche à modéliser :
\textit{des domaines aux caractéristiques différentes nécessitent des
stratégies de scoring différentes}.

\section{V1 vs V2 : impact de la pertinence sémantique}

L'un des choix techniques importants de ce travail est d'utiliser la
pertinence sémantique (V2, embeddings) plutôt que la pertinence lexicale
(V1, mots-clés) dans le scorer. Le tableau~\ref{tab:v1_v2} compare
les deux approches sur plusieurs domaines représentatifs.

\begin{table}[htbp]
  \centering
  \caption{Comparaison V1 (lexicale) vs V2 (sémantique) en NDCG@15.}
  \label{tab:v1_v2}
  \rowcolors{2}{lightgray}{white}
  \begin{tabular}{lccc}
    \toprule
    \textbf{Domaine} & \textbf{Optuna V1} & \textbf{Optuna V2} & \textbf{$\Delta$(V2-V1)} \\
    \midrule
    knowledge\_graph\_embedding & 0.599 & 0.602 & $+0.003$ \\
    explainable\_ai & 0.206 & 0.206 & $\pm 0.000$ \\
    computational\_biology & 0.547 & 0.566 & $+0.020$ \\
    text\_summarization & 0.263 & 0.681 & $+0.418$ \\
    quantitative\_finance & 0.099 & 0.226 & $+0.127$ \\
    \bottomrule
  \end{tabular}
\end{table}

La pertinence sémantique apporte un gain significatif pour les domaines
à vocabulaire spécialisé (biologie computationnelle, finance quantitative,
text summarization), où la correspondance lexicale stricte est insuffisante.
Pour les domaines à terminologie plus standard (knowledge graph, XAI),
les deux approches sont comparables. Ces résultats justifient le choix
de V2 comme version par défaut dans le système.

\section{Implémentation}

\subsection{Module \texttt{automl\_scorer.py}}

Le scorer AutoML est implémenté dans \texttt{src/scoring/automl\_scorer.py}.
La fonction principale \texttt{run\_optuna\_study()} orchestre :
\begin{enumerate}
  \item Le chargement de l'oracle depuis \texttt{data/oracle/\{domain\}/} ;
  \item La création d'une étude Optuna avec échantillonneur TPE et
    pruner MedianPruner ;
  \item L'optimisation de la fonction objectif NDCG@15 sur 150 essais ;
  \item La sauvegarde des pondérations optimales dans
    \texttt{data/weights/\{topic\}\_weights.json}.
\end{enumerate}

Les poids sont chargés automatiquement par \texttt{QualityCritic}
lors de l'exécution du pipeline, via la fonction
\texttt{load\_weights\_for\_topic(topic)}.

\subsection{Accélération par sous-scores précalculés}

Afin d'accélérer l'optimisation, les sous-scores ($s_f(a)$ pour chaque
article $a$ et feature $f$) sont calculés une seule fois avant la boucle
d'optimisation et mis en cache. L'évaluation d'un vecteur de pondérations
$\mathbf{w}$ se réduit alors à une multiplication matrice-vecteur suivie
du calcul de NDCG@15, ce qui prend moins d'une milliseconde. Avec 150
essais, l'optimisation complète d'un domaine nécessite moins de 30 secondes.
